% Chapter 22: The Kadison--Singer Problem
\let\LocalMacro\relax

\chapter{The Kadison--Singer Problem}

\section{The Problem}

\subsection{Informal}
In the 1950s, two mathematicians asked a question that sounds like a riddle: suppose you have a huge collection of vectors in space, each pointing in a different direction. Can you always split them into a few groups, each of which behaves almost as nicely as a set of mutually perpendicular axes? This question, which seemed purely abstract, turned out to have deep consequences for how we understand quantum measurements, signal processing, and even the structure of data.

\subsection{Simplified Version}
The problem is equivalent to the Paving Conjecture: given an infinite-dimensional diagonal operator $D$ with bounded entries, can the basis vectors be partitioned into finitely many blocks such that each block is nearly orthogonal? This was verified numerically for many examples before the proof, and Weaver had conjectured a finite-dimensional reformulation (Weaver's conjecture $KS_2$) that was already checked for small matrices. The breakthrough came from showing that Weaver's conjecture holds in general via the ``few moments'' method.

\subsection{Formal}
\begin{conjecture}[Paving Conjecture]
For every $\epsilon > 0$, there exists an integer $r$ such that any orthogonal projection $P$ on $\mathbb{C}^n$ with diagonal entries $P_{ii} \leq \epsilon$ can be $r$-paved. That is, the index set $\{1, \ldots, n\}$ can be partitioned into $r$ subsets $S_1, \ldots, S_r$ such that for each $j$, the submatrix $P_{S_j}$ satisfies:
\begin{equation}
\|P_{S_j}\| \leq 1 - \eta
\end{equation}
for some $\eta > 0$ depending only on $\epsilon$.
\end{conjecture}

This conjecture implies that one can decompose a large projection into smaller parts whose operator norms are strictly bounded away from 1.

\begin{historicalbox}
The Kadison--Singer problem remained unsolved for over 50 years. Its difficulty stemmed from its connection to diverse areas of mathematics, including C*-algebras, frame theory in signal processing, graph theory, and computer science.
\end{historicalbox}

\subsection{Why It Matters}
The Kadison--Singer problem sits at the crossroads of quantum physics, signal processing, and pure mathematics. A positive answer means that incomplete quantum measurements can always be completed consistently. This has implications for how we understand measurement in quantum mechanics and how we process incomplete data in engineering.

\subsection{What Was Known}
The problem was posed in 1959 and was known to be equivalent to several other important conjectures in mathematics, including the Paving Conjecture in operator algebras and the Feichtinger conjecture in harmonic analysis. By the 2000s, partial results had been obtained using techniques from Banach space theory, but the full problem remained open for over 50 years.

\section{Mathematical Theory}


\subsection{Prerequisites}
This chapter assumes familiarity with:
\begin{itemize}
\item Linear algebra (inner products, orthogonal projections, eigenvalues)
\item Operator theory and Hilbert spaces (bounded operators, pure states)
\item Basic graph theory (adjacency matrices, bipartite graphs)
\end{itemize}

\subsection{Frames and Discretization}
In signal processing, a \emph{frame} is a generalization of an orthonormal basis. A family of vectors $\{\phi_i\}_{i \in I}$ in a Hilbert space $H$ is a frame if there exist constants $A, B > 0$ such that:
\begin{equation}
A \|x\|^2 \leq \sum_{i \in I} |\langle x, \phi_i \rangle|^2 \leq B \|x\|^2
\end{equation}
for all $x \in H$. The Kadison--Singer problem is equivalent to the statement that any frame can be partitioned into a finite number of subsets, each of which is a frame with controlled bounds.

\subsection{Bipartite Ramanujan Graphs}
A $d$-regular graph is \emph{Ramanujan} if its adjacency matrix eigenvalues (excluding $\pm d$) have absolute value at most $2\sqrt{d-1}$. Ramanujan graphs are optimal expanders. In 2013, Marcus, Spielman, and Srivastava used the Kadison--Singer problem to prove the existence of bipartite Ramanujan graphs of all degrees.

\section{The Solution}

In 2013, Adam Marcus, Daniel Spielman, and Nikhil Srivastava solved the Kadison--Singer problem \cite{marcus2015a, marcus2015b}:

\begin{theorem}[Marcus--Spielman--Srivastava, 2013]
The Kadison--Singer Paving Conjecture is true. Specifically, for any vectors $v_1, \ldots, v_m$ in $\mathbb{C}^d$ such that $\sum_{i=1}^m v_i v_i^T = I$ and $\|v_i\|^2 \leq \epsilon$, there exists a partition of $\{1, \ldots, m\}$ into two sets $S_1, S_2$ such that:
\begin{equation}
\left\|\sum_{i \in S_j} v_i v_i^T\right\| \leq \frac{(1 + \sqrt{\epsilon})^2}{2}
\end{equation}
for $j=1,2$.
\end{theorem}

Their proof used a revolutionary new technique:

\begin{enumerate}
\item \textbf{Interlacing Families of Polynomials}: They showed that the roots of certain families of polynomials interlace, which guarantees the existence of a partition where the maximum root of the characteristic polynomial is small.
\item \textbf{Expected Characteristic Polynomials}: Instead of finding a partition directly, they studied the average characteristic polynomial over all random partitions, showing its roots are bounded by the roots of a associated univariate polynomial.
\item \textbf{Barrier Functions}: They used multivariate barrier functions to track and bound the roots of these polynomials under rank-one updates.
\end{enumerate}

\begin{highlightbox}
The Marcus--Spielman--Srivastava proof settled Kadison--Singer by converting an operator algebra problem into a question about the roots of polynomials. This "method of interlacing families" has since become a major tool in combinatorics and computer science.
\end{highlightbox}

\section{Impact}
\begin{itemize}
\item \textbf{Operator Algebras}: Resolved the 1959 Kadison--Singer C*-algebra conjecture.
\item \textbf{Signal Processing}: Proved that every frame can be paved, ensuring stable partitions in transmission codes.
\item \textbf{Computer Science}: Provided a deterministic construction of bipartite Ramanujan graphs of all degrees, solving a long-standing open problem in network design.
\end{itemize}

\section{Exercises}

\begin{exercise}[Basic]
Show that the union of two frames in $\mathbb{C}^d$ is a frame. If $A_1, B_1$ and $A_2, B_2$ are the frame bounds for the two frames, what are the bounds for the union?
\end{exercise}

\begin{exercise}[Intermediate]
Explain the relation between the Kadison--Singer problem and the construction of Ramanujan regular graphs.
\end{exercise}

\begin{exercise}[Challenge]
Prove that any orthogonal projection $P$ on $\mathbb{C}^n$ with diagonal entries $P_{ii} \leq \epsilon$ can be partitioned into two parts whose submatrices have norm bounded away from 1.
\end{exercise}

\begin{exercise}[Computational]
Write a Python/SageMath script to explore a concept from this chapter. See the appendix for starter code.
\end{exercise}

\section{Timeline}

\begin{tabularx}{\linewidth}{lX}
\toprule
\textbf{Year} & \textbf{Event} \\ \midrule
1959 & Kadison and Singer formulate the extension problem \cite{kadison1959} \\ 
1979 & Anderson reformulates the conjecture as the Paving Conjecture \cite{anderson1979} \\ 
1980s--2000s & Equivalent conjectures found in signal processing (Feichtinger conjecture) \\ 
2013 & Marcus, Spielman, and Srivastava resolve the conjecture using interlacing polynomials \cite{marcus2015a, marcus2015b} \\ 
2015 & Annals of Mathematics publishes the two-part proof \\ \bottomrule
\end{tabularx}

\section{Citations}

\begin{enumerate}
\item Kadison, R. V., \& Singer, I. M. (1959). ``Extensions of pure states.'' American Journal of Mathematics, 81(2), 383--400.
\item Anderson, J. (1979). ``Extreme points in secondary algebra.'' Journal of Functional Analysis.
\item Marcus, A., Spielman, D. A., \& Srivastava, N. (2015). ``Interlacing families I: Bipartite Ramanujan graphs of all degrees.'' Annals of Mathematics, 182(1), 307--325.
\item Marcus, A., Spielman, D. A., \& Srivastava, N. (2015). ``Interlacing families II: Mixed characteristic polynomials and the Kadison--Singer problem.'' Annals of Mathematics, 182(1), 327--350.
\end{enumerate}
